%%%%%%%%%%%%%%%%%%%%%%%%%%%%%%%%%%%%%%%%%%%%%%%%%%%%%%%%%%%%%%%%%%%%%%%%
% Plantilla TFG/TFM
% Universidad de Murcia. Facultad de Informática
%%%%%%%%%%%%%%%%%%%%%%%%%%%%%%%%%%%%%%%%%%%%%%%%%%%%%%%%%%%%%%%%%%%%%%%%


\chapter*{Extended Abstract}

Este extended abstract debe explicar de manera clara y resumida de qué trata el proyecto, por qué se eligió este tema y cuál es el objetivo principal del trabajo. Debe describir los métodos, herramientas o enfoques principales utilizados durante el desarrollo, destacar los resultados más relevantes obtenidos y resumir las conclusiones principales. También se puede indicar la importancia del trabajo y posibles aplicaciones o líneas futuras de desarrollo.  

Se recomienda que este texto esté redactado de forma que cualquier lector pueda entender el contenido y alcance del trabajo sin necesidad de leer el TFG completo.  

El extended abstract debe estar escrito en inglés y tener 2000 palabras.